\documentclass[12pt,a4paper]{article}
\usepackage{../../../myconfig}

\begin{document}

% =========================================================================
% BÀI TẬP ĐIỂM DANH: PHƯƠNG PHÁP TỌA ĐỘ HÓA / HÌNH HỌC KHÔNG GIAN (10 CÂU)
% =========================================================================
\titlebox{Chuyên đề: Oxyz}{Phương pháp tọa độ hóa}{Oxyz: Phương pháp tọa độ hóa}

\vspace{-0.25cm}

% Câu 1 (Câu 1 - Dạng 2 - Chuyên đề 4)
\questionLinesManual{0}{1}{%
Cho hình lập phương $ABCD.A'B'C'D'$ có cạnh bằng $a$. Khoảng cách giữa hai đường thẳng $BD$ và $A'C'$ bằng
\begin{tasks}(4)
    \task $\dfrac{\sqrt{3}a}{2}$.
    \task $\sqrt{2}a$.
    \task $\sqrt{3}a$.
    \task $a$.
\end{tasks}
}{}
\drawdashlinesfull

\vspace{0.15cm}

% Câu 2 (Câu 2 - Dạng 2 - Chuyên đề 4)
\questionLinesManual{0}{2}{%
Cho hình chóp $S.ABC$ có đáy là tam giác vuông tại $A$, $AB = 2a$, $AC = 4a$, $SA$ vuông góc với mặt phẳng đáy và $SA = a$. Gọi $M$ là trung điểm của $AB$. Khoảng cách giữa hai đường thẳng $SM$ và $BC$ bằng
\begin{tasks}(4)
    \task $\dfrac{2a}{3}$.
    \task $\dfrac{\sqrt{6}a}{3}$.
    \task $\dfrac{\sqrt{3}a}{3}$.
    \task $\dfrac{a}{2}$.
\end{tasks}
}{}
\drawdashlinesfull

\vspace{0.15cm}

% Câu 3 (Câu 4 - Dạng 2 - Chuyên đề 4)
\questionLinesManual{0}{3}{%
Cho hình chóp $S.ABC$ có đáy $ABC$ là tam giác vuông cân tại $A$, $AB = a$, $SA$ vuông góc với mặt phẳng đáy và $SA = a\sqrt{3}$. Gọi $M$ là trung điểm của $BC$. Khoảng cách giữa hai đường thẳng $AC$ và $SM$ bằng
\begin{tasks}(4)
    \task $\dfrac{a\sqrt{2}}{2}$.
    \task $\dfrac{a\sqrt{39}}{13}$.
    \task $\dfrac{a}{2}$.
    \task $\dfrac{a\sqrt{21}}{7}$.
\end{tasks}
}{}
\drawdashlinesfull

\vspace{0.15cm}

% Câu 4 (Câu 5 - Dạng 1 - Chuyên đề 4)
\questionLinesManual{0}{4}{%
Cho hình chóp $S.ABC$ có đáy là tam giác vuông đỉnh $B$, $AB = a$, $SA$ vuông góc với mặt phẳng đáy và $SA = 2a$. Khoảng cách từ điểm $A$ đến mặt phẳng $(SBC)$ bằng
\begin{tasks}(4)
    \task $\dfrac{2\sqrt{5}a}{5}$.
    \task $\dfrac{\sqrt{5}a}{3}$.
    \task $\dfrac{2\sqrt{2}a}{3}$.
    \task $\dfrac{\sqrt{5}a}{5}$.
\end{tasks}
}{}
\drawdashlinesfull

\vspace{0.15cm}

% Câu 5 (Câu 12 - Dạng 1 - Chuyên đề 4)
\questionLinesManual{0}{5}{%
Cho tứ diện đều $ABCD$ có cạnh bằng $a$. Tính khoảng cách từ đỉnh $A$ đến mặt phẳng $(BCD)$.
\begin{tasks}(4)
    \task $\dfrac{a\sqrt{6}}{2}$.
    \task $\dfrac{a\sqrt{6}}{3}$.
    \task $\dfrac{3a}{2}$.
    \task $2a$.
\end{tasks}
}{}
\drawdashlinesfull

\vspace{0.15cm}

% Câu 6 (Câu 1 - Dạng 1 - Chuyên đề 3)
\questionLinesManual{0}{6}{%
Cho tứ diện $OABC$ có $OA, OB, OC$ đôi một vuông góc với nhau và $OA = OB = OC$. Gọi $M$ là trung điểm của cạnh $BC$. Góc giữa hai đường thẳng $OM$ và $AB$ bằng
\begin{tasks}(4)
    \task $45^\circ$.
    \task $90^\circ$.
    \task $30^\circ$.
    \task $60^\circ$.
\end{tasks}
}{}
\drawdashlinesfull

\vspace{0.15cm}

% Câu 7 (Câu 3 - Dạng 1 - Chuyên đề 3)
\questionLinesManual{0}{7}{%
Cho hình hộp chữ nhật $ABCD.A'B'C'D'$, biết đáy $ABCD$ là hình vuông. Tính góc giữa hai đường thẳng $A'C$ và $BD$.
\begin{tasks}(4)
    \task $90^\circ$.
    \task $30^\circ$.
    \task $60^\circ$.
    \task $45^\circ$.
\end{tasks}
}{}
\drawdashlinesfull

\vspace{0.15cm}

% Câu 8 (Câu 2 - Dạng 2 - Chuyên đề 3)
\questionLinesManual{0}{8}{%
Cho hình chóp $S.ABC$ có $SA$ vuông góc với mặt phẳng $(ABC)$, $SA = a\sqrt{2}$, tam giác $ABC$ vuông cân tại $B$ và $AC = 2a$. Góc giữa đường thẳng $SB$ và mặt phẳng $(ABC)$ bằng
\begin{tasks}(4)
    \task $30^\circ$.
    \task $45^\circ$.
    \task $60^\circ$.
    \task $90^\circ$.
\end{tasks}
}{}
\drawdashlinesfull

\vspace{0.15cm}

% Câu 9 (Câu 3 - Dạng 3 - Chuyên đề 3)
\questionLinesManual{0}{9}{%
Cho hình chóp $S.ABC$ có $SA$ vuông góc với mặt phẳng $(ABC)$, $SA = \dfrac{a\sqrt{3}}{2}$, tam giác $ABC$ đều cạnh bằng $a$. Góc tạo bởi giữa hai mặt phẳng $(SBC)$ và $(ABC)$ bằng
\begin{tasks}(4)
    \task $90^\circ$.
    \task $30^\circ$.
    \task $45^\circ$.
    \task $60^\circ$.
\end{tasks}
}{}
\drawdashlinesfull

\vspace{0.15cm}

% Câu 10 (NÂNG CAO - ẨN CHIỀU CAO h: Câu 28 - Dạng 2 - Chuyên đề 3)
\questionLinesManual{0}{10}{%
Cho hình chóp $S.ABCD$ có đáy $ABCD$ là hình vuông cạnh $a$, cạnh bên $SA$ vuông góc với mặt phẳng đáy và $SA = h$ ($h > 0$). Biết góc giữa hai mặt phẳng $(SBC)$ và $(SCD)$ bằng $60^\circ$. Giá trị của $h$ bằng
\begin{tasks}(4)
    \task $h = 2a$.
    \task $h = a$.
    \task $h = \dfrac{3a}{2}$.
    \task $h = \dfrac{a}{2}$.
\end{tasks}
}{}
\drawdashlinesfull


% =========================================================================
% ĐÁP ÁN BÀI TẬP ĐIỂM DANH (10 CÂU TRÍCH TỪ TÀI LIỆU)
% =========================================================================
% 01. D | 02. A | 03. B | 04. A | 05. B
% 06. D | 07. A | 08. B | 09. C | 10. B

\end{document}